\chapter{实时与离线混合检索系统的需求分析与设计}
\section{系统总体概述} 
本章节首先从总体层面描述 Search Agent 系统的定位、功能边界与端到端工作流，为后续各模块的需求与设计奠定基础。Search Agent 的核心目标是在 MCP协调下，作为“提示（prompt）→检索/爬取→解析清洗→索引入库→结果返回”这一端到端流水线的执行与调度者。系统需要同时对接两类检索后端：一类是实时在线检索（以 bin search 为代表），具备热点定时检索与对检索结果触发 Python 实时爬取的能力；另一类是离线垂域检索（如小红书、知乎），以批量同步的方式提供丰富垂域内容，支持基于倒排索引的文本检索与基于向量索引的语义检索。Search Agent 的职责既包括将用户自然语言 prompt 转化为可执行检索指令，也包括根据业务策略在实时/离线通道间路由、触发爬取、控制并发与超时、管理任务的生命周期及错误降级，并负责将抓取或检索得到的数据以消息方式下发至数据处理流水线完成清洗与入库。
\begin{figure}[H]
\centering
\includegraphics[
width=1.\textwidth,keepaspectratio]{picc/系统介绍.drawio_1.pdf}
\caption{系统总体概述图}
\label{fig:系统介绍}
\end{figure}
图~\ref{fig:系统介绍}是系统的总体概述工作流，图中说明了系统的概述流程。图~\ref{fig:系统详细介绍}是系统的详细介绍图。结合两张图可以梳理清楚整个系统：用户通过上游系统发起带有 prompt 的检索请求，MCP 将请求路由至 Search Agent；Search Agent 先调用模型服务对 prompt 做语义解析与关键词提炼，产生结构化检索指令与元信息（意图、是否需要展示爬取详情、置信度等）；根据策略决定调用实时通道或离线通道或同时并行调用两者；若实时通道命中且解析出的关键词/元信息指示需要展示详情，则触发实时爬取任务并对爬取任务做速率/超时/优先级控制；检索或爬取产生的原始内容通过消息队列进入数据处理模块，进行实时解析、初步清洗、二次清洗与质量评估；合格数据被入库到倒排索引（文本检索）或向量索引（语义检索），以支撑后续检索请求；Search Agent 将召回候选进行融合与 TopK 合并，然后将候选送重排层做低延迟精排并返回最终结果，同时支持异步回调以在后续爬取或补充数据入库后更新展示。
\begin{figure}[H]
\centering
\includegraphics[
width=1.\textwidth,keepaspectratio]{picc/系统介绍.drawio_2.pdf}
\caption{系统详细分解图}
\label{fig:系统详细介绍}
\end{figure}
系统的关键非功能目标包括：端到端延迟在业务可接受范围内（对交互类请求 P95 延迟例如 < 300ms，P99 < 1s，为触发爬取的场景允许异步补全），高吞吐（，高可用（SLA 99.9\% 以上）、可扩展、可观测（丰富监控与告警）、以及合规安全。总体架构应采用分层、模块化设计以便拆分职责、便于单元化部署与灰度发布，同时在关键路径上采用缓存、异步化与降级策略以保证稳定性与响应性。
\section{系统需求分析}
\subsection{系统用例分析}
\begin{figure}[H]
\centering
\includegraphics[
width=1.\textwidth,keepaspectratio]{picc/系统用例图.pdf}
\caption{系统整体用例图}
\label{fig:用例图}
\end{figure}
图~\ref{fig:用例图}是本系统的总体用例图。本系统围绕智能检索与内容获取构建完整链路，其总体需求涵盖查询理解、任务调度、在线获取、离线召回以及数据处理与索引构建五个核心模块。首先，在Prompt解析与Query生成阶段，需要对用户输入进行语义解析与意图识别，提取关键实体、约束条件与上下文信息，并生成结构化查询表达，以支撑后续多通道检索。MCP协调与参数传递模块则负责统一调度各子系统，将解析后的查询映射为标准化参数，并根据策略将请求分发至不同检索通道，支持动态路由、策略控制以及多源结果融合的基础能力，同时保证参数传递的一致性与可追踪性。在实时在线检索与爬取触发方面，系统需具备低延迟，当现有索引无法满足需求时，自动触发爬虫任务进行补充抓取，实现对新鲜内容的覆盖，且需要控制抓取频率与资源消耗，保证系统稳定性。离线垂域检索模块侧重于高质量内容的深度召回，结合文本检索与向量检索提升语义匹配能力，依赖高质量索引与领域数据沉淀，强调召回的准确性与内容价值。最后，索引入库与数据处理模块承担数据清洗、去重、结构化处理与索引构建任务，是支撑整个系统性能与效果的基础，需保证数据处理流程标准化、可扩展，并支持增量更新与索引重建，以持续优化检索效果与系统整体表现。
\subsection{Prompt 解析与 Query 生成}
\begin{table}[H]
\centering
\caption{Prompt 解析与 Query 生成用例描述}
\label{tab:nonperiodic-usecase}
% 局部调整行高以增加留白（不要太紧凑）
{\renewcommand{\arraystretch}{1.}
\begin{tabular}{p{3cm} p{11cm}}
\hline
名称 & Prompt 解析与 Query 生成 \\ \hline
参与者 & 用户、前端应用、MCP、Search Agent、Prompt Parser、实时检索适配器、离线检索适配器（倒排/向量）、爬虫调度器、日志/审计服务 \\ \hline
触发条件 & 用户在前端提交自然语言检索请求或查询指令，请求被 MCP 接收并转发到 Search Agent。 \\ \hline
前置条件 & MCP 与 Search Agent 服务正常；Prompt Parser 的在线推理池可用并通过鉴权；实时与离线检索后端可访问；爬虫池处于允许状态（依据白名单与配额）。请求已附带必要上下文。 \\ \hline
后置条件 & 生成标准化的 QueryPackage 并记录解析元数据；根据策略路由至实时或离线通道（或并行）；当满足展示详情条件时触发爬取任务并写入任务队列；将解析与路由决策写入审计日志；将初步检索候选或占位结果返回给 MCP 供上游或前端展示，必要时通过异步回调补全详情。 \\ \hline
正常流程 &
MCP接收请求并鉴权，携带trace、上下文与优先级下发任务 -> Search Agent对prompt做规范化预处理并调用Prompt Parser轻量语义解析 -> Prompt Parser返回关键词、意图、展示指示、重写query、置信度与模型版本，规则引擎进行合规与安全校验并在必要时覆写或标记人工复核 -> Search Agent根据解析结果与全局策略（配额、白/黑名单、系统负载）路由请求（离线/实时/并行、向量优先或触发爬取）。\\
\hline
\end{tabular}
}
\end{table}
Prompt 解析模块是 Search Agent 的前置入口，其职责不仅是从自然语言中抽取关键词，更要识别用户意图、拆分多意图句子、生成可供各检索通道消费的结构化检索请求，并对是否触发爬取、检索范围与过滤条件给出明确建议。技术实现需融合规则与模型：规则层用于处理确定性逻辑（敏感词、白名单/黑名单、固定命令），模型层采用轻量化 encoder 或 sequence-to-sequence 模型做语义分词、实体识别与 query 重写。

输出必须包含明确的字段：关键词列表、扩展查询建议、意图标签、触发爬取布尔值与置信度分数，以便 MCP 与后端执行精确路由。响应时延是关键约束，交互路径要求 prompt 解析尽可能在几十到几百毫秒完成，故工程上需部署蒸馏模型、推理缓存与请求批处理策略；对于复杂解析可走异步增强路径，由离线更强模型定期校准解析器并提供回流修正。解析结果需可追溯（记录模型版本、解析耗时、置信度），并支持灰度发布与回滚，避免模型或规则迭代导致大范围质量回退。针对短文本与口语化输入，应加强同义词扩展、拼写容错与实体链接能力，以提升倒排与向量检索的召回覆盖。

\subsection{MCP 协调与参数传递} 
MCP 是系统的控制面，负责接收来自上游的请求、下发任务至 Search Agent 或爬虫池，并在异步任务完成后执行回调通知。对 MCP 的需求包括对同步与异步两类请求模式的统一支持、对任务状态生命周期的管理、以及对优先级与配额的全局调度能力。在同步场景，MCP 要求在预设延迟预算内等待并返回融合后的召回结果；在触发爬取并需异步补全的场景，MCP 需为任务分配唯一 task\_id 并提供任务状态查询、超时处理与回调 URL，以便前端或上游系统根据状态决定是否展示占位信息或发起主动刷新。参数传递方面，MCP 必须支持强类型接口，透明传递 trace\_id、模型版本、超时阈值、priority 等元信息，同时对任务的重试与幂等要求提供机制。

MCP 还承担资源配额控制，按域名/来源对爬取速率与并发上限进行限制，并在系统压力上升时提供全局熔断与降级策略。安全上，MCP 应实现鉴权与审计，记录每次参数变更与任务触发，确保满足合规审计要求。为提升可用性，MCP 需要水平扩展与 HA 设计，并将任务状态与元数据持久化以便故障恢复与历史回放。

表~\ref{tab:MCP 协调与参数传递用例描述}是MCP 协调与参数传递用例描述，MCP相当于是一个“负载均衡”的组件，为用户和Agent之间进行牵线搭桥，减少因为链路复杂带来的额外开销和维护成本。
\begin{table}[H]
\centering
\caption{MCP 协调与参数传递用例描述}
\label{tab:MCP 协调与参数传递用例描述}
% 局部调整行高以增加留白（不要太紧凑）
{\renewcommand{\arraystretch}{1.}
\begin{tabular}{p{3cm} p{11cm}}
\hline
名称 & MCP 协调与参数传递 \\ \hline
参与者 & MCP（协调平台）、前端/上游调用方、Search Agent、Prompt Parser、实时检索适配器、离线检索适配器（倒排/向量）、爬虫调度器、配额/白灰黑名单服务、消息队列、审计与监控服务 \\ \hline
触发条件 & 上游系统或用户通过前端提交检索请求到 MCP，或 Search Agent/其他组件发起需要 MCP 协调的异步任务。 \\ \hline
前置条件 & MCP 与各下游组件鉴权与网络连通正常；配置中心中存在任务协议（schema）、默认超时和重试策略、配额配置与回调约定；请求包含必要元数据。 \\ \hline
后置条件 & 请求被合法路由并携带完整参数透传到目标组件；任务状态在 MCP 中登记并可查询；异步任务注册了回调或事件订阅；所有关键参数变更与路由决策记录审计日志与监控埋点以支持回放与问责。 \\ \hline
正常流程 &
MCP接收请求并鉴权校验、限流，生成task\_id/trace\_id及超时与优先级 -> 将请求及元数据透传给Search Agent并记录路由，同步走等待路径，异步登记任务 -> 对需多下游调用的场景按预定义序列或并行策略拆分子任务并下发，调用均携带trace与鉴权与超时上下文 -> 与配额/白灰黑名单服务校验敏感操作权限，配额不足或合规禁止时拒绝或降级并记录计量 -> 同步路径聚合下游返回或超时后降级响应上游，异步路径维护任务状态机并在完成 -> 发生异常时执行重试/退避/降级/回滚策略，记录故障与决策写入审计日志并触发告警 -> 全流程埋点关键指标并汇总日志/trace到监控与回放系统以支持线上定位與策略优化。\\
\hline
\end{tabular}
}
\end{table}

\subsection{实时在线检索（bin search）与爬取触发} 
实时在线检索通道以低延迟与高新鲜度为核心目标，承担热点定时检索与对即时请求的快速响应。该通道需支持定期的热点扫描，将热点词推入检索队列并维持短时 cache 以应对突发流量。同时，Search Agent 必须在检索后根据 prompt 的解析结果与检索 snippet 的质量判断是否触发 Python 爬取以获得页面详情或补充信息。

触发条件既要保证用户体验（提高展示丰富度），又要控制系统与目标站点的压力，故需要策略性阈值、白名单/黑名单与实时配额。触发后的爬取由爬虫池异步执行，Search Agent 在主请求路径上应采取占位返还与异步补全策略：先返回检索候选或简短摘要以满足首屏响应需求，再在爬取完成后通过 MCP 推送更新以补全详情。

\begin{table}[H]
\centering
\caption{实时在线检索（bin search）与爬取触发模块用例描述}
\label{tab:nonperiodic-usecase}
% 局部调整行高以增加留白（不要太紧凑）
{\renewcommand{\arraystretch}{1.}
\begin{tabular}{p{3cm} p{11cm}}
\hline
名称 & 实时候选召回与爬取触发（bin search 路径） \\ \hline
参与者 & Search Agent、MCP、Prompt Parser、实时检索适配器（bin search client）、爬虫调度器、爬虫池、配额与白灰黑名单服务、监控与审计服务 \\ \hline
触发条件 & Search Agent 接收到来自 MCP 的已解析 QueryPackage，并判断需在实时通道进行低延迟候选召回。 \\ \hline
前置条件 & 实时检索适配器与后端 bin search 服务在 SLA 内可用；爬虫调度器及爬虫池处于允许接收任务的状态；配置中心含最新的域名白/黑名单与域级配额；Search Agent 与各适配器间鉴权与网络连通性正常；请求携带 trace\_id 和剩余超时预算。 \\ \hline
后置条件 & 实时检索适配器返回候选集合及 snippet；在满足触发条件且合规/配额允许下，爬取任务成功下发至爬虫调度器并入队；Search Agent 将初步候选用于首屏展示并订阅爬取回调以异步补全；所有路由与触发决策记录入审计日志并埋点上报监控。若爬取未被允许或失败，则按降级策略返回可替代展示。 \\ \hline
正常流程 &
MCP 下发带 QueryPackage到 Search Agent，Agent 检查剩余超时与优先级 -> 在短超时窗并行调用 bin search 适配器获取候选 -> 对候选做快速质量评估：snippet 密度、来源可信度（白/灰/黑名单、历史分）、与“需展示详情”匹配度，并读取域抓取可用性与域级配额 -> 若需抓取且域名允许，向配额服务申请并在获批后构造幂等爬取任务入队并记录审计-> 爬虫遵守严格超时/重试与退避策略，解析成功将结构化结果及 parse\_confidence、extraction\_errors 写入数据处理队列，失败记录原因并告警 -> 全流程埋点爬取率、配额使用、成功率、耗时与降级等指标，达到阈值时触发自动限流或人工干预。\\
\hline
\end{tabular}
}
\end{table}
实时检索通道对延迟与吞吐的双重约束要求在工程实现上做好缓存优先、并行调用倒排与向量后端、以及在高并发时进行部分降级（如返回缓存或减少重排深度）。此外，触发爬取前务必执行合规检查（robots、版权）与来源可信度评估，避免对受限制资源进行自动化抓取。

实时爬取模块需要在严格 SLA 下快速抓取目标页面并抽取结构化信息，以支撑展示或后续入库。该模块应在下载、渲染、解析三个阶段设定明确的超时阈值，并在发生超时时采取退避或降级策略。初步抓取优先采用轻量 HTTP 请求获取静态内容；若页面依赖大量 JS 渲染且需要完整 DOM 才能提取目标信息，则在探测到渲染必要性时切换到 headless 浏览器执行渲染，但应限制渲染并发与时间以控制成本。解析流程采用规则与模型融合的方法，先用基于规则的 readability 类库快速抽取正文，再以轻量 NER/IE 模型或正则规则进行字段抽取与标准化。解析结果需附带质量元数据（parse\_confidence、extraction\_errors、source\_score），以便后续清洗模块决定是否进行二次清洗或人工审核。为保证幂等性与避免重复，爬取任务在下发前应做 URL 规范化与指纹检查，任务执行应支持断点续传与重试上限。系统监控需覆盖成功率、超时率、平均/分位延迟、解析错误分布与队列堆积情况，并在阈值触达时触发自动降级（如暂停新的爬取任务或延长超时）以保护整体稳定性。安全与合规性包括对目标站点的请求频率控制、敏感信息过滤与爬取日志审计。

\subsection{离线垂域检索（文本检索/向量检索）}
离线检索通道面向小红书、知乎等社区垂域数据，强调语义丰富性与内容深度。离线数据通过定期批量采集或授权数据源同步进入 ETL 管道，经过清洗、去重、标签化与向量化后，分别构建倒排索引与向量索引以支持文本检索与语义检索。由于离线通道对实时性要求相对宽松，重点在于保证索引质量与语义一致性：embedding 模型版本的管理、预处理规则的统一、以及向量索引构建参数都需严格控制并记录以支持索引回滚或重建。离线数据的增量更新策略应兼顾新鲜度与构建成本，常见做法是在索引侧采用 delta 索引或周期性 bulk 重建并通过 alias 切换实现无缝上线。离线通道还需提供用于训练与评估的标注数据集，支持重排模型与召回融合的离线评估，进而为线上 Search Agent 的决策提供基准。质量控制包括对来源可信度、互动指标（点赞、评论）和文本完整性的评估，并将质量分作为索引字段供检索与重排模型使用。

此外，离线检索通道还需要构建完善的数据血缘与版本管理体系，对每一批数据的来源、处理流程以及索引构建过程进行全链路记录，以便在出现效果波动时能够快速定位问题并进行回溯分析。在特征工程层面，可以引入多粒度语义表示，如段落级与文档级向量，以提升长文本检索效果。同时，通过构建多路召回机制，将关键词匹配结果与语义召回结果进行融合，进一步提升召回覆盖率与相关性。在评估环节，应结合点击率、停留时长等行为数据构建离线指标体系，并与人工标注结果进行对齐，以形成稳定的评估基准。针对不同垂域场景，还可以设计差异化标签体系与领域词表，从而增强模型对专业语境的理解能力，最终提升整体检索质量与用户体验
\begin{table}[H]
\centering
\caption{离线垂域检索（文本检索/向量检索）模块用例描述}
\label{tab:nonperiodic-usecase}
% 局部调整行高以增加留白（不要太紧凑）
{\renewcommand{\arraystretch}{1.}
\begin{tabular}{p{3cm} p{11cm}}
\hline
名称 & 离线垂域检索（文本检索与向量检索）用例 \\ \hline
参与者 & Search Agent、MCP（请求协调）、Prompt Parser、离线倒排索引服务、向量检索服务、索引管理/元数据服务（Index Manager）、特征/模型服务、监控与审计服务 \\ \hline
触发条件 & Prompt Parser 或 Search Agent 判断该请求应使用离线通道（例如需高语义覆盖、垂域社区来源 XHS/知乎 优先或用户指定离线来源），由 MCP 将请求下发至 Search Agent 并路由到离线检索模块。 \\ \hline
前置条件 & 离线倒排索引与向量索引已按计划构建并对外提供查询接口；索引版本（index\_version）和 embedding\_model\_version 可查询并与请求兼容；检索服务健康且鉴权通过；请求包含 trace\_id、期望候选数与超时预算。 \\ \hline
后置条件 & 返回整合后的候选集合（包含来源、score、index\_version 与质量元数据）；将低质量或无结果的请求及其原始上下文记录到回放/告警系统以便后续补采或索引改进；检索全程元数据写入审计日志以便离线分析。 \\ \hline
正常流程 &
Search Agent接收QueryPackage并选离线通道，计算剩余超时与优先级 -> 并行发起关键词文本检索与向量检索，指定index/embedding版本与检索参数 -> 倒排与向量服务在超时内返回候选，超时则回退到另一路或使用缓存 -> 对候选按doc\_id/内容指纹去重并归一化/融合得分形成候选池 -> 延迟允许时运行轻量LTR并对top-N启用精排，否则返回LTR结果 -> 以统一格式回传MCP，记录检索指标并提交回放样本以支持离线分析与索引优化。 \\
\hline
\end{tabular}
}
\end{table}
\subsection{索引入库与数据处理模块} 

数据处理模块负责将来自实时检索、爬取与离线导入的原始内容转化为可搜索、合规与高质量的索引文档。实时管道要求低延迟处理：对爬取产物先做轻量化解析与去噪（HTML 去壳、编码修复、简单实体抽取），并快速计算质量分与指纹后入队进行二次清洗。
\begin{table}[H] \centering \caption{索引入库与数据处理模块用例描述} \label{tab:nonperiodic-usecase} 
% 局部调整行高以增加留白（不要太紧凑） 
{\renewcommand{\arraystretch}{1.} 
\begin{tabular}{p{3cm} p{11cm}} \hline 
名称 & 索引入库与数据处理 \\ \hline 参与者 & 爬虫/检索后端（产出原始内容）、消息队列、数据处理模块、索引构建服务（倒排索引构建器、向量化与向量索引器）、索引管理/元数据服务、Search Agent、监控与审计服务 \\ \hline 触发条件 & 爬虫或检索后端将原始抓取内容或离线导入批次写入消息队列或对象存储，触发数据处理流水线进行实时解析/批量清洗并最终入库构建倒排或向量索引。 \\ \hline 前置条件 & 数据处理服务与索引服务可用并通过鉴权；消息队列或对象存储中存在待处理的原始数据或批次；索引构建参数在索引管理服务中可查；系统具有足够的写入配额与资源配额以执行入库任务。 \\ \hline 后置条件 & 清洗后的结构化文档通过幂等写入入倒排索引或向量索引；索引元数据更新并持久化；生成入库日志与审计记录以便回放与溯源；对低质量或异常文档标记并触发人工复核或重处理任务。 \\ \hline 正常流程 &
数据来源写入消息队列/对象存储并登记元数据-> 实时/批量消费者拉取并快速解析，生成结构化记录 -> 结构化记录二次清洗，质量评估产出，低质或疑似敏感标记人工复核或降权 -> 通过阈值的记录按路径入库：文本写入倒排索引，需向量化的送Embedding并写入向量索引，写入需幂等 -> 索引构建完成后更新索引元数据并通过alias原子切换使新数据可见，记录耗时与成功率用于监控与告警 -> 入库失败或质量不达标写入回放并触发重处理，全部处理版本与关键元数据记录审计以支持可追溯与合规。 \\ \hline \end{tabular} } \end{table}

二次清洗在离线或异步流程中进行更复杂的处理，包括语义一致性校验、重复检测、实体规范化、时间与作者字段标准化、以及使用模型判别的内容可信度评估。质量评估体系需结合规则和机器学习：规则覆盖格式校验与明确违禁项，模型用于语义相似性、内容完整性和敏感信息检测。处理流程应记录完整的数据血缘与处理版本，以支持审计与问题回溯。清洗后的数据需按照不同索引需求生成不同格式的入库包：倒排索引侧注重字段标准化与分词质量，向量索引侧注重 embedding 的稳定性与元数据一致性。对低质量或可疑条目，系统应支持自动降权、暂缓入库或人工复核机制，并将复核结果用于改进自动化质量判定模型。

索引入库阶段要求高可用、可扩展且具备可观测的写入体系。对于倒排索引，需要设计批量导入与近实时刷新的混合策略，保证离线数据批量导入的压缩与查询性能，同时允许对重要实时数据采取 delta flush 以满足短期可搜索性。对于向量索引，需考虑向量化模型版本、构建参数与向量插入的代价：大规模重建常采用离线批处理，实时新增向量可采用内存热表或小型 HNSW 以实现短期可检索，随后做定期合并。入库过程应由消息队列驱动并支持幂等写入（使用唯一 ID 去重），并在发生写入失败或索引切换时提供原子化切换与回滚流程。写入还需记录索引元数据，以便后续查询兼容性与回溯。为保证服务连续性，索引写入要具备速率控制和后端退压策略，并监控写入延迟、错误率与索引大小增长，结合 MCM 做容量与维护窗口规划。

\subsection{非功能性需求分析} 
体系的非功能性需求涵盖性能、可用性、可扩展性、安全性与可运维性。在性能方面，需对不同场景设定明确的服务水平协议（SLA）。对于交互类请求，应设定P95延迟目标，例如小于300ms，并且确保P99延迟小于1秒。对于触发爬取的场景，允许异步补全处理，但需要保证补全的时延在可接受的范围内，通常为数秒级。吞吐量方面，要求系统具备每秒并发检索的能力以及最大并发爬取的数目，并确保在高负载下依然能够保持稳定的响应能力和数据处理速度。此外，针对大规模数据处理，系统应具备动态扩展能力，以应对未来可能的流量增长。

在可用性方面，要求系统具备冗余机制、自动重试机制和快速回滚能力，确保系统在任何故障情况下能迅速恢复。整体的系统可用性应保持在99.9\%以上，以保证用户的高可用需求。此外，系统还应具有全面的故障转移和灾难恢复机制，确保在发生严重故障时能够快速切换到备用系统，避免对业务造成较大影响。

可扩展性方面，系统需要支持组件的水平扩展能力，确保各个模块（如检索后端、爬虫池、模型推理服务和处理消费者）都可以根据负载自动进行扩展。同时，要求系统具备高效的分片副本策略，保证数据的分布和副本能够支持大规模并发的查询和数据处理需求。在一致性方面，要求系统的索引入库必须提供最终一致性保证，并且要尽可能减少索引数据的延迟，以保证用户的搜索体验不受影响。

安全性方面，系统必须具备完善的鉴权与授权机制，确保用户的身份和访问权限得到严格管理。此外，必须实施全面的审计日志功能，记录系统的所有关键操作，以便追溯。对于敏感信息，系统需进行脱敏处理，同时满足相关的合规性要求，包括版权和隐私的保护。同时，系统需要对爬取行为进行实时合法性校验，确保数据来源的合法性和合规性。

可运维性要求方面，系统应具备完善的监控和告警机制，能够实时监控系统的各项指标，并在异常情况发生时自动触发告警。同时，系统应支持自动化运维工具，如滚动升级、灰度发布和容量预警，以确保系统在高负载和复杂环境下依然能够稳定运行。为了保障系统的高可靠性，必须建立全面的测试策略，包括单元测试、集成测试、压力测试、回放测试和混沌演练等。通过这些测试手段，能够有效评估系统在不同故障场景下的恢复能力和韧性，从而提高系统的整体稳定性。
\section{系统总体设计}
\subsection{系统逻辑视图}
系统逻辑视图从功能与角色角度描述各子系统之间的职责分工与交互关系，便于理解 Search Agent 在端到端流水线中的定位与控制边界。总体上，系统由五类核心逻辑组件组成：请求协调层、决策与路由层、模型推理与解析层、检索与爬取后端、以及数据处理与索引层。每类组件在逻辑视图中承担明确职责：MCP 负责接收上游请求并维持任务生命周期、配额与回调；Search Agent 负责将 prompt 转为可执行检索指令、在实时/离线通道间做路由决策、触发爬取并维护任务并发与超时策略；Prompt Parser 为 Search Agent 提供语义分词、意图识别、关键词提炼与置信度评估；实时检索（bin search）用于低延迟候选召回并可在条件满足时触发爬虫获取详情；爬虫池执行实时抓取并将原始结果发往数据处理模块；离线索引服务（面向 XHS、知乎等源）提供基于倒排与向量索引的深度召回；数据处理层执行实时解析、去噪、质量评估并将合格数据落入适配的索引库。

图\ref{fig:系统逻辑视图}展示了一个典型的搜索代理系统的自上而下的分层逻辑视图。在该系统中，服务层、后端与中间件层以及存储层各自承担了不同的功能和任务，共同支撑整个搜索流程的高效运转。
\begin{figure}[H]
\centering
\includegraphics[
width=1.\textwidth,keepaspectratio]{picc/逻辑视图.pdf}
\caption{系统逻辑视图}
\label{fig:系统逻辑视图}
\end{figure}


首先，服务层负责接收来自客户端的请求，进行预处理和路由决策。请求通过前端进行验证和追踪，随后由Prompt Parser进行规范化并提取关键词和意图信息。路由器根据策略将请求分发至合适的检索后端，分为实时搜索和离线搜索。如果需要，系统还会请求Embedding Service进行文本到向量的转换。爬取管理器负责决策是否触发新的爬取任务，并通过爬取调度器将任务队列传递至爬虫执行器，最终将抓取结果写入消息队列并触发异步补全。

后端与中间件层主要负责搜索请求的具体处理，包括实时和离线的索引查询、嵌入服务的向量生成以及缓存管理等。该层中的消息队列提供了松耦合的数据传递通道，支持系统的高效异步操作。存储层包括对象存储、索引存储和元数据数据库，用于存储原始抓取内容、倒排索引、向量索引及相关的元数据。通过这些组件的协作，系统能够高效地处理用户请求，进行数据抓取与处理，并提供精确的搜索结果。
在逻辑交互上，整个请求生命周期从 MCP 收到带有 metadata的 prompt 开始。MCP 将请求分发给 Search Agent，Search Agent 首先调用 Prompt Parser 进行语义解析，得到结构化输出。基于解析结果与全局策略，Search Agent 决定是先调用离线通道、实时通道、还是同时并行调用两者。离线通道通常直接调用索引服务的检验API，实时通道则先向 bin search 发起请求以快速获取候选，若 snippet/候选不足且解析标识需展示详情则触发爬虫池。本模块设计上强调“先快后补”的理念：在触发爬取的场景下，Search Agent 会先返回可用的候选或占位摘要以保证交互响应，同时在后台等待爬取完成并在完成后通过 MCP 的回调或事件通知上游更新展示。

召回融合逻辑在逻辑视图中作为独立职责被强调：来自倒排索引与向量索引以及实时检索/爬取的候选需在融合层做去重、score normalization 与 TopK 合并，然后送入在线重排层。重排层采用分层重排策略：首先用轻量 feature-based 排序器做快速排序，再对 top-N 使用更强的 cross-encoder或摘要生成模型对展示内容做最终加工。数据处理模块在此逻辑链路中既承担向索引的落库职责，也负责对爬取内容做实时清洗与质量判别，将质量元数据反馈给召回与重排模块以影响候选权重。

为了支持系统演进与运维，逻辑视图中还包括配置与控制子系统、以及策略管理子系统。这些横向能力通过统一的控制面下发策略并确保 Search Agent 在运行时具有可控的行为边界。图示描述建议绘制一张逻辑架构图：中心为 Search Agent，左侧为 MCP 与上游请求，右侧并列展示实时检索与离线索引，下方为爬虫池与数据处理，顶部为模型推理服务与配置；箭头标注主要数据流（prompt→parse→recall→crawl→process→index→response）与控制流。
\subsection{系统过程视图}
系统过程视图从时序与流程角度刻画关键业务场景的执行步骤，便于分析延迟、容错与故障恢复逻辑。以三类典型流程为主：同步检索请求流程、触发爬取的交互流程、与离线数据入库。
\begin{figure}[H]
\centering
\includegraphics[
width=1.\textwidth,keepaspectratio]{picc/过程视图1.pdf}
\caption{数据爬取清洗解析入库时序图}
\label{fig:zhouqishixian}
\end{figure}
对于同步检索流程，当用户发起带 prompt 的请求时，MCP 作为入口进行身份与配额校验并将请求派发到 Search Agent；Search Agent 立即调用 Prompt Parser 进行语义解析，并行向模型服务发起轻量推理请求，解析返回关键词/意图并带有置信度；Search Agent 并行调用实时检索和/或离线索引服务，在收到候选后执行快速的 TopK 合并，并将候选送入轻量重排器计算最终排序；若在预设延迟窗口内无法获得充分候选，则 Search Agent 可返回占位与部分结果，并在后台继续等待其他通道的补充。整个同步流程需在严格 SLA 内完成，过程中 Search Agent 维护 trace id 并将关键耗时点埋点以便 Prometheus/Grafana 展示。

触发爬取的交互流程比同步流程多出实时爬取与异步回调环节。流程开始如上，但当 Prompt Parser 或候选质量判断出需要页面级详情时，Search Agent 构造爬取任务并将任务通过消息队列或 MCP 下发到爬虫池，同时仍以最快速度返回检索候选或摘要以保证首屏体验。爬虫池接到任务后执行下载，在下载与解析阶段执行严格的超时控制；解析完成后将结构化结果写入消息队列或直接调用数据处理服务的入队接口。数据处理模块对爬取产物做实时解析、清洗、质量评估并在合格后将条目落库到对应索引；当入库或解析完成时，系统通过 MCP 的回调机制将更新推送给上游或前端，从而实现异步补全与用户界面的数据刷新。整个链路要求幂等、去重与失败补救。

离线数据入库与索引重建流程侧重于批量处理与一致性保障。数据采集端将原始数据批次写入对象存储并触发 ETL 作业；数据处理层使用批处理框架进行大规模清洗、向量化，并生成索引构建任务。索引构建采用隔离式构建，在临时索引环境构建完毕后使用别名/atomic switch 原子切换到新索引，减少对线上查询的影响。对于需要较短可搜索延迟的新增数据，系统会采用 delta 索引策略，先将新数据写入小规模近实时索引，并在合适时机合并回主索引
\subsection{系统开发视图}
开发视图描述了搜索系统内部的核心组件及其交互关系，展示了系统的逻辑层次和模块划分。应用层中的各个服务组件通过容器化部署在Kubernetes集群中，负责处理来自客户端的请求，包括请求验证、语义解析、路由决策、检索调用等。服务之间通过gRPC、HTTP和消息队列进行通信，确保数据流转和任务调度。检索服务通过实时和离线索引服务提供搜索功能，而爬虫调度和执行模块则负责抓取和解析新的网页数据。数据处理和索引模块对爬取的数据进行清洗、质量评估和索引写入，确保数据的完整性和可搜索性。该架构具备高扩展性和灵活性，可以快速适应变化的需求。

图~\ref{fig:系统开发视图}展示了搜索系统的整体概述架构，通过不同的组件层次展示了系统的整体结构与模块间的交互流程。整个系统架构分为多个层次，包括网关/API层、应用层（Search Agent）、后端与爬取层、数据处理与索引层、基础设施层以及跨切面库层。系统首先由MCP Gateway负责接收请求并进行身份验证、配额校验以及路由决策。请求会通过IQueryAPI接口转发到应用层的相关模块进行处理。

在应用层，Front-End模块会验证请求并进行初步的追踪，随后由Prompt Parser模块对用户输入的prompt进行规范化处理，并提取关键词和意图信息。此时，Router模块根据请求类型做出路由决策，将请求发送至不同的后端服务进行查询。如果请求需要实时数据，Recall Adapter模块会调用实时搜索引擎（如Bins Cluster）；如果请求涉及离线数据，模块则会向离线索引系统（Offline Index Cluster）发起查询。若在应用过程中，需要新的数据源或页面，Crawl Manager模块会判断并构造爬取任务，通过消息队列提交到爬虫池，由爬虫执行器（Crawler）进行数据抓取与解析，随后通过消息队列将结果发送到数据处理模块。

数据处理与索引层负责对爬取或查询的结果进行解析、清洗和处理。Parser \& Cleaner模块会对抓取的数据进行解析并清洗，同时生成文档指纹。接着，Indexer模块将处理后的数据写入倒排索引和向量索引，并将索引数据持久化到Index Storage中，确保系统能够高效检索和查询。这一过程中的数据也会被存储到对象存储中，确保数据的可靠性和持久性。
\begin{figure}[H]
\centering
\includegraphics[
width=1.\textwidth,height=0.6\textheight]{picc/开发视图.pdf}
\caption{系统开发视图}
\label{fig:系统开发视图}
\end{figure}


在基础设施层，Message Queue提供了系统间异步通信的能力，确保模块之间能够高效且松耦合地进行数据交换。Object Storage则用于存储原始的网页数据和其他静态资源。系统还通过Monitoring模块实时监控各项指标，记录日志和操作，帮助进行故障排查与性能优化。

最后，跨切面库层包括一些提供全局功能的库，如AuthLib负责处理系统的权限和安全策略，而TraceLib则用于请求的追踪和系统配置管理，确保系统可以进行全面的性能分析和调优。系统中的每个模块和组件都与其他模块通过定义良好的接口和消息队列进行交互，从而实现高效的数据处理和查询能力，确保系统在满足性能和可靠性要求的同时，还能够灵活地进行扩展和维护。


\subsection{系统物理视图}
物理视图展示了搜索系统的部署架构，涵盖从客户端到基础设施层的完整部署流程。系统的请求首先通过客户端发起，经由API网关进行身份验证、限流、路由决策和追踪处理后转发到Kubernetes集群中的应用层。应用层的多个微服务组件负责请求验证、Prompt解析、路由决策、召回适配和爬取任务管理等功能。后端的检索集群、爬虫执行集群和数据存储系统则确保系统能够高效地处理查询请求和抓取新数据。在基础设施服务层，消息队列、对象存储、索引存储和元数据数据库保证了系统的数据流转和持久化。监控和运维系统提供了实时的系统监控、日志管理和故障追踪，确保系统的高可用性和稳定性。

系统物理视图描述实际部署拓扑、资源分配、网络与安全边界及高可用配置。推荐将系统部署在 Kubernetes 集群上以获得弹性扩缩容与运维统一化管理。关键运行单元包括：Search Agent 服务（stateless，可水平扩展）、Prompt Parser 、爬虫池、消息队列集群、数据处理，以及索引服务。

图~\ref{fig:物理视图}为系统的物理部署图，该图展示了一个复杂的搜索系统的物理架构，包括从客户端到基础设施服务的完整部署流程。整个系统的请求首先由客户端发起，经过API网关进行鉴权、限流、路由以及追踪等处理。网关将请求转发至Kubernetes集群中的核心微服务层，其中包括前端请求验证、Prompt解析、路由决策、召回适配、爬取任务管理、候选融合等功能模块。

在Kubernetes集群中，微服务通过容器化的方式部署，每个服务模块承担不同的功能。前端服务负责请求验证并传播追踪信息，Prompt解析服务将用户的查询转化为规范化的关键词和意图，路由决策服务则基于实时或离线检索需求决定请求的处理方式。召回适配器通过并行调用实时检索集群和离线索引服务，从后端检索系统获取候选结果。这些候选结果会被融合、去重并精排，最终返回给API网关。

当请求需要爬取新数据时，路由服务会触发爬取任务的创建，爬取任务会提交给专门的爬虫调度器，爬虫池中的执行器根据任务要求抓取网页，解析内容并上传到对象存储中，同时将处理后的数据通过消息队列传递至数据处理模块。数据处理模块对爬取的数据进行解析、清洗、质量评估后将合格的结果写入倒排索引或向量索引中，最终存储在索引存储系统中。
\begin{figure}[H]
\centering
\includegraphics[
width=1.\textwidth,height=0.7\textheight]{picc/物理视图.pdf}
\caption{系统物理视图}
\label{fig:物理视图}
\end{figure}


基础设施服务层包括消息队列、对象存储、索引存储和元数据数据库，确保各个服务之间的高效通信、数据持久化与备份。消息队列用于异步传输任务和数据，支持系统的高并发与解耦。对象存储用于保存原始页面与快照，而索引存储则用于持久化搜索索引，元数据数据库记录任务和数据的相关信息。
此外，系统还提供了强大的可观测与运维支持，涵盖了Prometheus和Grafana用于监控、ELK和Loki用于日志集中管理、Jaeger和Zipkin用于分布式追踪。所有的操作都被监控并记录，及时触发警报以保障系统的稳定运行。CI/CD流程和镜像仓库则确保了系统的快速迭代与部署，通过自动化构建和测试流程，确保系统组件的持续更新和高效交付。整个系统通过模块化的部署和容器化管理，实现了高可扩展性、灵活性与高可用性。

\section{系统模块设计}
\subsection{系统层次与模块划分}
系统采用分层模块化设计，以保证职能清晰、可扩展与易运维。最外层为接入与协调层，由MCP承担请求接入、鉴权、全局流量与配额控制、任务状态管理与回调；其下为控制与决策层的Search Agent，负责prompt的接收、调用解析模块、根据策略路径到实时或离线通道、触发爬取并维护任务并发与超时；解析与模型层提供语义分词、意图识别、query重写与置信度评估，是Search Agent的语义中枢；召回层包含实时在线检索与离线检索，并以统一接口向上暴露候选集合；爬取层由独立的爬虫池和渲染池组成，负责下载、渲染与初步解析；数据处理层承接检索、二次清洗、质量评估与入库；索引层分为倒排索引服务与向量索引服务，分别承担文本与语义检索能力；最后是重排与展示层，负责多通道候选融合、TopK合并、低延迟重排及摘要生成以供前端呈现。各层通过明确的API或消息schema解耦，关键通信使用gRPC或Protobuf以保证高效与向后兼容，同时消息驱动的队列用于解耦异步任务。

模块划分以职责单一为原则，每个模块内部再按功能细化子模块。例如Search Agent内部分为路由引擎、策略引擎、任务管理器与监控埋点组件；Prompt Parser划分为轻量推理路径与精校验路径；数据处理层分为流式解析通道与批处理通道，分别承担实时清洗与离线深度处理。模块之间约定统一的元数据模型包含trace\_id、task\_id、source、model\_version、priority、timeout等，以实现端到端可追溯。为支持灰度与多版本并行，系统在模块层面设计了版本控制与特性开关，并在MCM平台中集成发布审批、灰度流量控制与回滚机制。监控与告警覆盖所有关键模块的SLI，包括延迟、错误率、吞吐、队列堆积、爬取成功率、解析质量分布等，并设计了自动化的降级与回退策略以降低事故影响。

Prompt解析模块是将自然语言prompt转化为检索指令与任务元信息的核心。该模块设计为“多层混合架构”，上层为轻量化在线解析器，保证低延迟响应，面向交互路径提供快速的关键词抽取、意图判别与可控的queryrewrite；下层为高质量离线解析器，用于周期性或异步地对线上解析结果进行回流校正与策略更新。在线解析器由两部分构成：规则引擎与蒸馏模型。规则引擎处理确定性逻辑，如指令式prompt、安全或合规预检、白名单或黑名单命中，并对模型输出做约束。蒸馏模型为经过蒸馏或量化的encoder或seq2seq模型，侧重在短文本语义抽取，并通过缓存常见prompt解析结果与请求批处理来降低延迟与成本。
\subsection{Prompt 解析模块设计} 

\begin{figure}[H]
\centering
\includegraphics[
width=1\textwidth,keepaspectratio]{picc/prompt模块流程图(1).pdf}
\caption{Prompt 解析模块设计流程图}
\label{fig:zhouqishixian}
\end{figure}
解析模块的输出为标准化的QueryPackage，包括关键词列表、扩展查询（同义词、别名）、意图标签、触发爬取标志、置信度score以及上下文元数据。该输出必须具备可解释性，模型需要返回中间的attention或score以便调试，并且在低置信度场景下触发fallback逻辑，自动降级为规则解析或请求用户进一步澄清。为了保证解析的一致性，模块维护预处理规范，如分词、去噪、停用词策略、截断规则等，以及后处理策略，如同义词扩展表、实体对齐规则等。这些配置通过配置中心下发，支持在运行过程中进行灰度调整，以满足不同业务需求和灵活应对数据变化。
\begin{figure}[H]
\centering
\includegraphics[
width=1.\textwidth,keepaspectratio]{picc/prompt模块设计类图.pdf}
\caption{Prompt 解析模块设计设计类图}
\label{fig:zhouqishixian}
\end{figure}


在部署与运维方面，解析模块需要支持多版本并行、动态切换以及A/B实验。低延迟路径会部署在靠近Search Agent的推理池中，使用ONNX（开放神经网络交换格式）或ORT（ONNX Runtime）等工具实现高效推理；高质量模型则部署在GPU池中，并以异步方式服务于离线校准或复杂请求的处理。在这种架构下，模型的响应速度和精确度能够得到最优化，同时支持对不同版本的模型进行实时切换，以便进行迭代和性能调优。为了保证系统的安全与合规性，解析模块会对敏感信息进行实时识别与脱敏处理。日志层面，所有的prompt都会进行哈希化存储，这不仅有助于减少敏感信息的暴露，还能满足审计和合规性需求，确保系统的透明性和合规性。

性能监控是解析模块的另一个重要环节，模块通过监控解析延迟分位数、模型置信度分布、fallback率以及解析引起的下游爬取触发率等关键指标，持续评估模型的表现并对其进行优化。这些监控指标为模型的优化与规则阈值调整提供了数据支持，能够帮助团队识别系统瓶颈和改进点。此外，实时的数据反馈也能够帮助调整模型的参数，从而确保系统在不同场景下都能保持较高的鲁棒性和效率。通过这些精细的监控机制，解析模块不仅提升了系统的响应速度和准确度，同时也加强了系统的稳定性和可靠性。

为了满足不断变化的业务需求，解析模块还需适应动态场景的变化。例如，在某些情况下，用户的请求可能会涉及到更多复杂的查询，或是涉及到不同领域的知识，这就要求解析模块具有较强的自适应能力。通过将推理池部署在不同的服务层级，并结合多版本和动态切换机制，解析模块能够灵活地应对不同层次的查询需求，不仅能提供高速的实时响应，还能支持高精度的复杂场景分析，满足业务的多样化需求。

\begin{figure}[H]
\centering
\includegraphics[
width=1.\textwidth,keepaspectratio]{picc/prompt模块时序图.pdf}
\caption{Prompt 解析模块设计时序图}
\label{fig:zhouqishixian}
\end{figure}

\subsection{实时检索与爬取模块设计}
实时检索与爬取模块承担两类任务：第一在极短时间内提供候选召回（bin search 适配器），第二在判定需要更详尽内容时触发 Python 爬取并完成实时解析。实时检索适配器设计为轻量客户端池，维持与 bin search 的连接池与并发控制，侧重在调用延迟与稳定性上做工程优化（并行请求、请求合并、重试与熔断）。热点定时检索作为后台任务维护高频词表并定期刷新结果到本地 cache，用于降低突发流量对后端的冲击。检索返回后，Search Agent 根据解析输出与 snippet 质量判断是否触发爬取。

\begin{figure}[H]
\centering
\includegraphics[
width=1\textwidth,keepaspectratio]{picc/爬取模块流程图(1).pdf}
\caption{实时检索与爬取模块流程图}
\label{fig:zhouqishixian}
\end{figure}

爬取模块为一个可弹性扩缩的爬虫池，内部分为下载层、渲染层与解析层。下载层优先进行轻量 HTTP 获取，只有在探测到页面依赖 JS 或关键数据缺失时才升级到 headless 渲染层（Playwright/Puppeteer）。为控制成本与遵守目标站点友好策略，渲染实例池采用资源隔离（独立 GPU/CPU 机器组或独立容器组）并对渲染并发设定域名级与全局级配额。解析层则负责正文抽取、结构化字段提取与初步质量打分（parse\_confidence），常用方法为规则（readability、XPath）与轻量模型（DOM 节点分类器、NER）结合。爬取任务以消息队列驱动，任务描述包括 url、render\_needed、max\_fetch\_time、priority 与 callback\_topic，任务执行需保证幂等与断点续传特性，并在完成后将解析结果以结构化消息写回数据处理队列。

超时与并发控制是实时爬取的核心保障。每个任务设定严格的全链路超时（下载超时、渲染超时、解析超时），并采用指数退避和有限重试策略。为避免爬取对目标站点造成压力，系统维护域名白/灰/黑名单及爬取策略（速率、深度限制），在检测到高失败率或被封禁风险时自动降级。安全与合规方面，爬取前需检查 robots.txt 规则并在必要时进行人工授权；对可能涉及敏感信息的页面，解析结果需自动触发脱敏或人工复核流程。监控维度包括爬取成功率、平均耗时、渲染占比、解析错误率与队列堆积，出现异常自动触发告警与限流策略。


\begin{figure}[htb]
\centering
\includegraphics[
width=0.8\textwidth,keepaspectratio]{picc/爬取模块设计类图.pdf}
\caption{实时检索与爬取模块设计类图}
\label{fig:zhouqishixian}
\end{figure}

\begin{figure}[htb]
\centering
\includegraphics[
width=1.\textwidth,keepaspectratio]{picc/爬取模块时序图.pdf}
\caption{实时检索与爬取模块时序图}
\label{fig:zhouqishixian}
\end{figure}

\subsection{离线检索与索引模块设计}
离线检索与索引模块面向垂域数据（XHS、知乎等）构建高质量的倒排索引与向量索引。整个离线管线由数据采集、批量清洗、向量化、索引构建与索引发布五个阶段组成。数据采集支持 API 拉取、第三方授权导出或批量文件导入，采集过程记录完整元数据以支持血缘追踪。批量清洗在批处理框架下完成，包含去重、HTML 去壳、字段标准化、语言检测与初步质量评分，清洗任务需产出可重放的中间结果以便回放重跑。

向量化阶段使用统一的 embedding 模型与预处理规范，embedding 训练，所有索引记录 embedding 模型版本以便于后续回滚或重建。向量索引构建使用 Faiss 等库，构建参数在离线调优后记录入库。倒排索引构建采用 Lucene，并在构建结束后通过别名切换（atomic alias switch）无缝替换旧索引以保证查询可用性。
\begin{figure}[htb]
\centering
\includegraphics[
width=0.8\textwidth,keepaspectratio]{picc/模块流程图(1).pdf}
\caption{离线检索与索引模块流程图}
\label{fig:zhouqishixian}
\end{figure}

\begin{figure}[htb]
\centering
\includegraphics[
width=1.1\textwidth,keepaspectratio]{picc/模块设计类图.pdf}
\caption{离线检索与索引模块设计类图}
\label{fig:zhouqishixian}
\end{figure}
增量更新策略采用delta 索引与周期性合并，实现对实时或近实时数据的高效处理。小批量写入的数据首先存储在 delta 索引（或热表）中，查询时系统并行访问主索引与 delta 索引，确保低延迟检索。达到预定阈值后，后台任务会自动将 delta 索引合并到主索引，并进行优化，以减少冗余数据并提高查询性能。为了保证索引的一致性和准确性，入库过程要求幂等写入（避免重复数据）和事务记录（包括写入日志和校验和），确保数据完整性。

索引的元数据管理涵盖版本、构建参数、文档计数与快照信息，帮助运维人员进行回滚与容量规划。系统通过实时监控查询延迟、写入延迟、索引增长以及段合并时间分布，确保索引性能的稳定性。此外，支持离线评估，索引模块提供快照回放接口，使得历史索引可切换到测试环境中进行回归验证或A/B 实验，进一步优化系统的稳定性和准确性。
\begin{figure}[H]
\centering
\includegraphics[
width=1.\textwidth,keepaspectratio]{picc/模块时序图.pdf}
\caption{离线检索与索引模块时序图}
\label{fig:zhouqishixian}
\end{figure}
\subsection{召回融合与重排模块设计}

\begin{figure}[H]
\centering
\includegraphics[
width=1\textwidth,keepaspectratio]{picc/模块4流程图(1).pdf}
\caption{召回融合与重排模块设计流程图}
\label{fig:zhouqishixian}
\end{figure}

召回融合与重排模块负责将多通道候选（实时倒排、离线倒排、向量检索、爬取结果）合并、去重并给出最终排名，是影响最终用户感知质量的关键环节。融合阶段首先对不同通道的原始分数进行归一化或校准，以消除评分尺度差异。去重逻辑在候选合并前后均执行：先按 doc\_id 简单去重，再对文本相似度或内容指纹进行近似去重以处理文本近似的重复条目。合并策略依据业务权重控制不同通道的候选份额，并在合并后做候选裁剪。
\begin{figure}[H]
\centering
\includegraphics[
width=1.\textwidth,keepaspectratio]{picc/模块4设计类图.pdf}
\caption{召回融合与重排模块设计设计类图}
\label{fig:zhouqishixian}
\end{figure} 

重排采用分层策略以兼顾质量与延迟：第一层为轻量 feature-based ranker（LightGBM/GBDT），输入特征包含原始检索分数、匹配特征（BM25、term overlap）、语义相似度（embedding cosine）、爬取/解析质量元数据、来源可信度与用户上下文特征；该层以毫秒级延迟返回排序结果。第二层（可选）为 cross-encoder 或深度排序模型，用于对 top-K（例如 top50）进行精确再排序，以提升最终展示相关性。为降低第二层延迟，可采用蒸馏技术将 cross-encoder 的知识迁移至轻量模型，或仅在高价值请求/付费流量上启用 cross-encoder。摘要生成与最终展示也可在重排后并行执行，摘要模型要受延迟预算约束，可异步返回更长摘要。
\begin{figure}[H]
\centering
\includegraphics[
width=1.\textwidth,keepaspectratio]{picc/模块4时序图.pdf}
\caption{召回融合与重排模块设计时序图}
\label{fig:zhouqishixian}
\end{figure}
在线部署则需考虑模型推理并发、特征预计算与缓存。为降低在线延迟，许多特征（如文档静态质量分、来源可信度、历史交互统计）可以预先离线计算并存储在特征仓库，实时请求只需取用。重排模块需监控关键指标（重排延迟、重排提升的 CTR/满意度、跨通道贡献度），并在 CI/CD 中对模型上线进行离线回归与在线小流量灰度验证。最终，融合与重排输出应携带足够的解释性信息（top features、source attribution）以便前端展示“结果为何被选中”的透明提示并支持审计。
\section{数据库与索引存储设计}
本节围绕系统中承担数据持久化、索引存储与检索支持的数据库与索引层进行系统性设计说明。设计目标是在保证检索质量与数据新鲜度的前提下，兼顾写入吞吐、查询延迟、可扩展性、可恢复性与合规安全。整体存储分为三类原始抓取/中间产物存储（对象存储为主）、结构化/元数据存储（关系型或 NoSQL 数据库）、以及检索索引存储（倒排文本索引与向量索引）。各类系统在逻辑上相互配合，通过消息队列与元数据引用完成端到端的数据流转与可追溯性。

原始数据与中间产物存储采用对象存储（S3/MinIO）作为单一真相（source of truth）。所有爬取产物、第三方导出文件、以及离线批处理中间文件（例如分片向量文件、原始 HTML 快照）都以时间和任务维度组织在对象存储中，并保留完整的元信息（task\_id、source、fetch\_time、content\_sha256）。对象存储的优势在于强持久性、成本可控及便于回放重跑。当出现解析模型或清洗逻辑问题时，可基于对象存储的原始快照重新执行流水线。为支持快速检索与短期回放，系统在对象存储之上还维护短期的 Redis/NoSQL 缓存，例如将近期抓取的正文或解析结果的嗅探摘要保存为缓存键，便于低延迟展示与幂等检查。


\begin{table}[H]
\centering
\caption{用户查询表：存储用户输入的查询信息}
\label{tab:user_queries}
{\renewcommand{\arraystretch}{1}
\begin{tabular}{|p{3cm}|p{3cm}|p{6cm}|}
\hline
\textbf{字段名} & \textbf{数据类型} & \textbf{描述} \\ \hline
query\_id & INT & 用户查询的唯一标识 \\ \hline
user\_id & INT & 用户的唯一标识 \\ \hline
query\_text & TEXT & 用户输入的查询文本 \\ \hline
query\_timestamp & DATETIME & 查询时间 \\ \hline
keywords & TEXT & 由模型提取出的关键词 \\ \hline
\end{tabular}
}
\end{table}
user\_queries 表是系统中存储用户查询信息的核心表格。它记录了每个用户提交的查询内容以及与之相关的其他信息。首先，query\_id 是每次查询的唯一标识，确保能够唯一地识别用户的每个查询。user\_id 则是标识执行查询的用户的唯一标识符，通常是通过用户注册或认证系统分配的，用于区分不同的用户。query\_text 存储用户输入的查询文本，即用户在搜索框中输入的具体内容。这些查询文本会被进一步传递给模型进行分词处理，提取出相关的关键词。

query\_timestamp 字段记录了查询的时间，这对于日志分析、用户行为分析等非常重要。通过时间戳，可以了解用户活跃的时间段，以及是否存在某些季节性或突发的查询需求。keywords 则是通过模型对用户查询文本进行处理后提取出来的关键词，关键词的提取是系统进行后续搜索和内容筛选的关键。这些关键词不仅影响到搜索引擎的查询结果，还决定了是否需要爬取外部数据或进行更复杂的内容展示。

这个表格的设计支持快速存储和检索用户的查询信息，为后续的数据处理、爬取和质量评估提供基础数据。同时，它为系统的日志记录、用户行为分析以及查询分析提供了宝贵的参考。
\begin{table}[H]
\centering
\caption{搜索结果表：存储搜索引擎返回的结果}
\label{tab:search_results}
{\renewcommand{\arraystretch}{1}
\begin{tabular}{|p{3cm}|p{3cm}|p{6cm}|}
\hline
\textbf{字段名} & \textbf{数据类型} & \textbf{描述} \\ \hline
result\_id & INT & 搜索结果的唯一标识 \\ \hline
query\_id & INT & 关联的用户查询ID \\ \hline
source\_type & VARCHAR(50) & 搜索源类型（如在线搜索、离线搜索） \\ \hline
content\_id & INT & 搜索到的内容的唯一标识 \\ \hline
content\_url & TEXT & 搜索结果的URL或路径 \\ \hline
content\_summary & TEXT & 搜索结果的摘要 \\ \hline
content\_timestamp & DATETIME & 内容抓取的时间 \\ \hline
status & VARCHAR(20) & 搜索状态（如成功、失败、爬取中） \\ \hline
\end{tabular}
}
\end{table}
search\_results 表用于存储从搜索引擎返回的搜索结果。每次用户提交查询后，系统会根据查询文本和提取的关键词调用搜索引擎（如实时在线搜索或离线搜索）进行检索，返回的结果将被记录在此表格中。

首先，result\_id 字段为每条搜索结果提供了一个唯一标识，确保每条记录能够被单独识别。query\_id 则与 user\_queries 表中的 query\_id 关联，标明当前搜索结果对应的是哪次查询。这使得我们能够追踪到每个用户的查询以及该查询的返回结果。

source\_type 字段指示该搜索结果来自哪个搜索源，可能的值包括“在线搜索”和“离线搜索”，这有助于区分实时数据和历史数据的来源。content\_id 用于关联 content\_library 表中的内容记录，指明该搜索结果具体是哪个内容。content\_url 存储了该搜索结果的具体链接或路径，可以是一个网页URL，也可以是一个本地存储路径。

content\_summary 字段提供了该搜索结果的简短摘要，帮助用户快速了解搜索结果的概要信息。content\_timestamp 记录了搜索结果的抓取时间，帮助系统追踪抓取的时间窗口，便于在内容更新时及时刷新结果。最后，status 字段记录了该搜索结果的处理状态，如“成功”、“失败”或“爬取中”，这些状态对于调试和监控爬虫任务非常重要。
\begin{table}[H]
\centering
\caption{爬取的数据表：存储从爬虫抓取的数据}
\label{tab:crawled_data}
{\renewcommand{\arraystretch}{1}
\begin{tabular}{|p{3cm}|p{3cm}|p{6cm}|}
\hline
\textbf{字段名} & \textbf{数据类型} & \textbf{描述} \\ \hline
crawled\_id & INT & 爬取数据的唯一标识 \\ \hline
content\_id & INT & 关联的内容ID \\ \hline
raw\_content & TEXT & 原始爬取的内容 \\ \hline
cleaned\_content & TEXT & 清洗后的内容 \\ \hline
crawl\_timestamp & DATETIME & 爬取时间 \\ \hline
status & VARCHAR(20) & 爬取状态（如成功、失败） \\ \hline
\end{tabular}
}
\end{table}
crawled\_data 表记录了通过爬虫从外部网站抓取的原始数据。随着在线搜索的进行，如果搜索引擎认为需要提供更详细的内容，系统会启动爬虫从指定的页面进行抓取。该表格的作用是存储这些爬取到的原始数据，并为后续的数据清洗和处理提供基础。

每条爬取数据都有一个唯一的标识符 crawled\_id，以便于管理和追踪。content\_id 用于与 content\_library 表中的内容记录关联，帮助我们识别哪些内容是通过爬虫获取的。raw\_content 存储了爬虫抓取到的原始内容，可能包含杂乱的信息、HTML标签等。cleaned\_content 则存储了清洗后的内容，去除冗余和无关信息，使其更加规范和有用。

crawl\_timestamp 字段记录了数据的爬取时间，这对于判断数据的新鲜度和时效性至关重要。爬虫的状态通过 status 字段进行标识，常见的状态包括“成功”和“失败”。爬虫状态的管理对于优化爬虫流程、排查错误及防止爬虫陷入死循环等问题非常有帮助。
\begin{table}[H]
\centering
\caption{内容库表：存储所有的内容，包括文本、图像等}
\label{tab:content_library}
{\renewcommand{\arraystretch}{1}
\begin{tabular}{|p{3cm}|p{3cm}|p{6cm}|}
\hline
\textbf{字段名} & \textbf{数据类型} & \textbf{描述} \\ \hline
content\_id & INT & 内容的唯一标识 \\ \hline
content\_type & VARCHAR(50) & 内容的类型（如文本、图像） \\ \hline
content\_data & TEXT & 内容数据 \\ \hline
timestamp & DATETIME & 内容录入时间 \\ \hline
status & VARCHAR(20) & 内容状态（如正常、已删除） \\ \hline
\end{tabular}
}
\end{table}
content\_library 表是整个系统中的核心内容存储库，所有的内容数据都会在这里进行归档。这个表格记录了所有通过搜索引擎、爬虫或其他方式获得的内容，无论是文本数据、图像还是其他格式的数据。该表为系统后续的搜索结果展示和数据分析提供了稳定的支持。

content\_id 是该表的主键，用于唯一标识每一条内容记录。content\_type 字段用于标记内容的类型，常见的类型包括文本、图像、视频等，这对于后续的内容处理非常重要。content\_data 则存储了具体的内容数据，可能是纯文本、图像的URL链接，或者是其他类型的数据。

timestamp 字段记录了该内容的录入时间，这对于内容的版本控制、时效性评估和更新非常重要。status 字段则指示内容的状态，如“正常”表示该内容有效，“已删除”表示该内容被移除，帮助系统在展示时能够过滤掉无效或过时的内容。
\begin{table}[H]
\centering
\caption{处理后的数据表：存储经过清洗和质量评估后的数据}
\label{tab:processed_data}
{\renewcommand{\arraystretch}{1}
\begin{tabular}{|p{3cm}|p{3cm}|p{6cm}|}
\hline
\textbf{字段名} & \textbf{数据类型} & \textbf{描述} \\ \hline
processed\_id & INT & 处理数据的唯一标识 \\ \hline
raw\_data\_id & INT & 关联的原始数据ID \\ \hline
cleaned\_data & TEXT & 清洗后的数据 \\ \hline
quality\_score & FLOAT & 数据质量分数 \\ \hline
timestamp & DATETIME & 处理时间 \\ \hline
status & VARCHAR(20) & 处理状态（如成功、失败） \\ \hline
\end{tabular}
}
\end{table}

processed\_data 表用于存储经过清洗和质量评估的数据，旨在提升数据的准确性和可靠性。数据清洗是从爬虫获取原始数据到最终展示内容的关键一步，目的是删除不必要的噪声、格式化数据并规范化内容。这个表的设计确保了清洗后的数据能够更好地用于后续的分析和展示。

processed\_id 字段是该表的主键，每条处理过的数据都有唯一的标识。raw\_data\_id 关联了原始数据的 ID，通过与 crawled\_data 表的关联，可以追溯每条处理数据的来源。cleaned\_data 存储了清洗后的数据，这些数据已被标准化并且去除了冗余信息，确保了数据的高效利用。

quality\_score 字段是一个浮动值，用于表示处理数据的质量评分。质量评估通常基于数据的完整性、准确性和时效性等多方面的标准，这对于系统的可靠性至关重要。timestamp 字段记录了数据处理的时间，便于管理和版本控制。status 字段则指示处理状态，如“成功”表示数据已处理完成，“失败”则表示处理过程中出现了问题。

\begin{figure}[H]
\centering
\includegraphics[
width=0.6\textwidth,keepaspectratio]{picc/er图.pdf}
\caption{关系型数据数据库设计er图}
\label{fig:zhouqishixian}
\end{figure}
结构化元数据与任务状态存储采用关系型数据库与 NoSQL结合的方式。核心元数据包括文档元信息表、爬取任务表、索引元数据表、质量评估表与审计日志表。WebContentDB 为检索条目的目录中心，记录 document\_id、canonical\_url、title、source、fetch\_time、processing\_version、quality\_score、index\_status等关键字段；设计时需为 document\_id 与 canonical\_url 建立唯一约束以保证去重幂等写入。爬取任务表记录 task\_id、url、priority、attempts、last\_status、worker\_id、start\_time、end\_time 与 error\_reason，用于任务恢复、溯源与运维。索引元数据表记录每次索引构建的版本号、embedding\_model\_version、构建参数、segment\_info 与 snapshot\_location，以支持回滚与 A/B 检验。审计日志表记录对关键数据或配置的审批与变更轨迹，保证合规审计要求。

倒排索引与向量索引分别采用专门的索引系统来满足不同检索特性，确保系统在处理传统文本检索与基于语义的检索时都能具备高效性和准确性。倒排索引是传统信息检索的基础，适用于高频的短文本检索。它将文档中的每个词条（term）与包含该词的文档建立倒排映射。推荐使用成熟的分布式搜索引擎来实现倒排索引系统。这些引擎提供了丰富的查询表达能力，支持高级的查询语言与过滤功能，能有效满足用户复杂检索需求。此外，它们具备强大的分词与分析器定制功能，可以根据特定应用场景对中文或其他语言的分词进行优化，从而提高检索的精度。分片与副本机制则是分布式系统中不可或缺的组成部分，它们保障了索引的扩展性与容灾能力，在处理海量数据时能平衡负载并提供高可用性。

为了支持中文语义与短文本prompt的检索，倒排索引的mapping设计非常关键。首先，需要为不同字段使用差异化的分词器，尤其是在处理中文时，可以采用细粒度的分词策略，避免语义丢失。同时，为了提高短语匹配的能力，针对需要短语匹配的字段，索引设计时应保留position信息，这样不仅能支持高效的短语检索，也能保证文档的语义完整性。此外，设置term vectors（术语向量）可以使检索引擎对每个词汇的出现位置、频率和邻接关系有更加详细的记录，这对于后续的重排特征计算至关重要，能够显著提高检索结果的相关性与准确性。

在倒排索引的分片策略方面，建议采用逻辑哈希分配分片的方式，以便在大规模数据集上能够均匀分布查询负载并支持横向扩展。此外，为了保障高可用性和容灾能力，系统应配置副本机制，确保在某一分片出现故障时，其他副本能够承担查询请求，保证系统的正常运行。对于高写入场景，倒排索引的segment merge策略需要进行调优，以避免合并操作时产生的I/O瓶颈。在实际应用中，segment merge时的频繁I/O可能会导致系统出现性能瓶颈，特别是在实时更新场景下。为了减少这些影响，系统应优化合并策略，合理调整合并的阈值和频率，避免对系统吞吐产生过大压力。

针对实时通道的低延迟需求，倒排索引服务应启用near-real-time配置，以确保索引能够在较短时间内反映数据的更新。这意味着索引刷新过程应尽量减少对查询性能的影响。在此基础上，可以结合delta索引策略来进一步优化写入性能。Delta索引允许系统将新增的文档或数据记录先写入一个小规模的近实时索引，当达到一定条件后，再将这些增量数据合并到主索引中。这样不仅能减少索引刷新频率带来的合并压力，还能有效提升实时查询的性能。

向量索引则专门用于语义搜索，尤其是在处理复杂查询如短文本或非结构化数据时，能够显著提高检索结果的相关性。向量索引通常基于先进的向量化技术，如Word2Vec、BERT等模型生成的向量表示，并依赖于高效的向量检索引擎（如FAISS、Milvus等）来提供相似度搜索功能。与倒排索引相比，向量索引更侧重于通过语义相似度计算来完成检索，而不仅仅依赖于词汇的精确匹配。因此，向量索引的建库方式需要更多地考虑向量化的质量以及计算效率。为了提高检索性能，向量索引系统通常会采用高效的近似最近邻（ANN）搜索算法，通过构建K近邻图或其他数据结构来加速大规模向量空间中的相似度查询。

为了支持向量索引的高效存储与检索，系统必须对向量数据进行有效的索引构建。常见的向量索引方式包括基于树的数据结构（如KD树、Ball树）、基于量化的索引（如PQ、IVF）以及图搜索算法（如HNSW）。在建库过程中，系统需要根据实际场景的需求选择合适的索引策略。例如，对于实时性要求较高的场景，向量索引系统可以结合多层次索引机制，将高维向量空间划分为多个子空间，通过索引优化加速检索速度；而对于批量数据，系统可以采用静态批量索引方式，以保证高效的存储和查询性能。
\section{本章小结}
本章从模块设计延伸到数据库与索引存储的具体架构，系统性地回答了 Search Agent 在数据持久化、检索索引与可运维性方面的关键工程问题。模块划分明确了控制、解析、检索、爬取、处理与重排各自的职责边界；Prompt 解析模块在延迟与质量之间采用规则与模型的混合策略以实现可控可审计的输出；实时检索与爬取模块通过任务化、域名限速与严格超时保证低延迟路径的稳定性；离线索引模块通过批量处理与 delta 策略平衡新鲜度与压缩效率；数据处理模块建立了规则+模型的质量评估闭环，支持自动降级与人工复核；召回融合与重排模块采用分层设计兼顾效果与性能。数据库与索引存储设计部分则提出了面向工程实践的落地方案：对象存储作为原始快照、关系型/NoSQL 存储负责元数据与任务管理、倒排与向量索引分别承担文本与语义检索，并通过别名切换、备份快照与索引版本管理实现可回滚与灰度。整套设计在保证检索质量与新鲜度的同时，也兼顾了写入吞吐、查询延迟、扩展性与合规性。下一章将基于本章设计给出关键模块的实现细节、接口定义、伪代码示例与性能优化实验方案。